% Rationale Zahlen – bank/rationale-zahlen/ (zusammenbau v0.4, Heft)
\einheitenkopf[t5]{Rationale Zahlen}
\setcounter{aufgabe}{38}

% Anlauf (ohne Prüfkennung)
\begin{aufgabe}{Ordnen – Anlauf}
\begin{teile}
\teil Schreibe an jeden Strich seine Zahl. Abstand zwischen zwei Strichen: $1$.

\zahlenstrahl[xmin=-5,xmax=5,xstep=1]{0/0}
\teil Trage ein: $-4$, $3$ und $-8$.

\zahlenstrahl[xmin=-10,xmax=10,xstep=1]{}
\teil $-15$ oder $-9$ – welche ist kleiner? \leerfeld
\end{teile}
\end{aufgabe}

\begin{aufgabe}{Ordnen – Prüfungsaufgaben}
\begin{teile}
\teil Gib eine Zahl an, die größer ist als $-320$. \hfill (P10 2015 OS) \\ \leerfeld
\teil Ein Tauchroboter soll höher bleiben als $-85$ m. Nenne eine mögliche Position in m. \hfill (P10 2015 OS) \\ \leerfeld[m]
\teil Ordne aufsteigend: $-\frac{2}{5}$; $1{,}3$; $-0{,}43$; $0{,}9$ \hfill (P10 2014 OS) \\ \leerfeld $<$ \leerfeld $<$ \leerfeld $<$ \leerfeld
\teil Ordne aufsteigend: $-\frac{1}{4}$; $2{,}6$; $-0{,}257$; $0{,}5$ \hfill (P10 2014 OS) \\ \leerfeld $<$ \leerfeld $<$ \leerfeld $<$ \leerfeld
\end{teile}
\end{aufgabe}

% Anlauf (ohne Prüfkennung)
\begin{aufgabe}{Multiplizieren/Dividieren – Anlauf}
\begin{teile}
\teil $(-6) \cdot 7$ – Vorzeichen des Ergebnisses?\\ \kreuz{plus} \kreuz{minus}
\teil Berechne: $7 \cdot (-5)$ \leerfeld
\teil Berechne: $(-7) \cdot (-6)$ \leerfeld
\end{teile}
\end{aufgabe}

\begin{aufgabe}{Multiplizieren/Dividieren – Prüfungsaufgaben}
\begin{teile}
\teil Berechne ohne Taschenrechner den Wert von $\frac{a + b}{c}$ für $a = 15$, $b = -6$ und $c = -4$. \hfill (P10 2021 OS) \\ \leerfeld
\teil Berechne ohne Taschenrechner den Wert von $\frac{a + b}{c}$ für $a = 5$, $b = -2$ und $c = -4$. \hfill (P10 2021 OS) \\ \leerfeld
\teil Berechne den Wert von $(a + b) : c$ für $a = 3$, $b = -9$ und $c = -2$. \hfill (P10 2016 OS) \\ \leerfeld
\teil Berechne den Wert von $(a + b) : c$ für $a = 5$, $b = -19$ und $c = -7$. \hfill (P10 2016 OS) \\ \leerfeld
\teil Berechne den Wert von $3 \cdot (x - 5)$ für $x = -4$. \hfill (P10 2026 FOR) \\ \leerfeld
\teil Berechne den Wert von $6 \cdot (x - 2)$ für $x = -3$. \hfill (P10 2026 FOR) \\ \leerfeld
\teil Eine negative Zahl wird mit einer negativen Zahl multipliziert. Das Ergebnis ist … \hfill (P10 2015 OS)\\ \kreuz{immer positiv}\\ \kreuz{immer negativ}\\ \kreuz{Das kann man nicht entscheiden.}
\teil Eine negative Zahl wird durch eine positive Zahl geteilt. Das Ergebnis ist … \hfill (P10 2015 OS)\\ \kreuz{immer positiv}\\ \kreuz{immer negativ}\\ \kreuz{Das kann man nicht entscheiden.}
\end{teile}
\end{aufgabe}

% Anlauf (ohne Prüfkennung)
\begin{aufgabe}{Terme und Sachaufgaben – Anlauf}
\begin{teile}
\teil Auf Mias Konto sind $120$ €. Dann werden $45$ € abgebucht. Was ist der Startwert, was die Änderung? Startwert \leerfeld , Änderung \leerfeld
\teil Kontostand $-35$ €, Einzahlung $50$ € – neuer Kontostand? \leerfeld[€]
\teil Löse die Klammer auf und berechne: $15 - (8 - 17)$ \leerfeld
\end{teile}
\end{aufgabe}

\begin{aufgabe}{Terme und Sachaufgaben – Prüfungsaufgaben}
\begin{teile}
\teil Tierpark, Preise: Erwachsene $17$ €; Kinder von 4 bis 15 Jahren $8$ €; Kinder unter 4 Jahren frei; Familienkarte (2 Erwachsene und bis zu 3 Kinder) $46$ €; Großeltern-Enkel-Karte (2 Großeltern mit 1 oder 2 Enkeln, nur Montag bis Freitag) $30$ €. Am Mittwoch kommen Mutter, Vater, Oma, Opa und die Kinder Lina (3 Jahre), Paul (8) und Emma (9). Ermittle den günstigsten Gesamtpreis und vergleiche mit den anderen Möglichkeiten. \hfill (P10 2016 OS) \\ Günstigster Preis: \leerfeld[€]
\teil Freizeitbad, Preise: Erwachsene $15$ €; Kinder und Jugendliche von 6 bis 17 Jahren $7$ €; Kinder unter 6 Jahren frei; Familienkarte (2 Erwachsene und bis zu 2 Kinder) $36$ €; Gruppenkarte (5 Personen, nur Samstag und Sonntag) $38$ €. Am Samstag kommen zwei Eltern mit Ida (4 Jahre), Ben (7), Mara (15) und Tim (17). Ermittle den günstigsten Gesamtpreis und vergleiche mit den anderen Möglichkeiten. \hfill (P10 2016 OS) \\ Günstigster Preis: \leerfeld[€]
\teil 2024 kamen $48\,315$ Menschen zum Stadtfest, das sind $3\,870$ mehr als 2023. Wie viele waren es 2023? \hfill (P10 2017 OS) \\ \leerfeld Menschen
\teil Heute zeigt das Thermometer $-3$ °C, das sind $5$ Grad mehr als gestern. Wie kalt war es gestern? \hfill (P10 2017 OS) \\ \leerfeld °C
\end{teile}
\end{aufgabe}
